% ===================================================================
%  TABLEAU N° 14 — La Rue. — Les Commerçants.
%  Feuillets 78 a 83 du fac-simile = folios 76 a 81.
%  Correspondance : folio imprime = numero de feuillet - 2.
%
%  Les cles « %%K » sont les MEMES que dans texto/io : c'est par elles
%  que la page de lecture met les deux textes en regard. Le francais ne
%  connait aucun des dix intertitres que l'ido insere (« Exploro al
%  butiki. », « La kukiferio. », « Marcho-defilo. », etc.) : Guignon
%  laisse le recit courir sans coupure, la ou Rochelle le decoupe en
%  scenes. Ces cles-la restent donc sans vis-a-vis.
%
%  ATTENTION — LE FAC-SIMILE FRANCAIS EST UNE PRISE DE VUE, non un
%  passage au scanner : contraste et echelle varient d'une page a
%  l'autre. Les renvois chiffres ont ete recoupes avec
%  texto/io/23-tabelo-14.tex chaque fois qu'un chiffre etait ambigu
%  (ex. « colporteur » (69), « marchand d'habits » (86)) : la meme
%  legende numerote la meme gravure dans les deux livrets.
% ===================================================================

% --- Feuillet 78 (folio 76 : ouverture de tableau) ------------------
\begin{VUpage}[78]{76}
%%K t14-apar-1 apar
\VUsaut{4.0mm}
\VUtitre{15.0pt}{60}{TABLEAU N\textsuperscript{o} 14}
\VUsaut{3.0mm}

%%K t14-tit-1 sub
\VUtitre{11.4pt}{0}{La Rue. --- Les Commerçants.}
\VUsaut{3.0mm}

%%K t14-01-1 p
\VUblancAlinea
1. --- Mon ami Jean, qui habite la campagne, est venu\nl
passer trois jours dans ma famille. Jusqu'à présent il\nl
n'avait pas encore eu l'occasion de voir une ville\nl
importante, aussi nous employons toutes nos journées\nl
à nous promener et je lui explique ce qu'il ne connaît\nl
pas.

%%K t14-02-1 p
\VUblancAlinea
2. --- Nous sommes allés tout d'abord visiter\nl
l'\VUgras{observatoire} \textsuperscript{(1)}, qui l'a beaucoup intéressé. En\nl
revenant par la \VUgras{rue} \textsuperscript{(3)} où se trouve l'\VUgras{établissement}\nl
\VUgras{de bains turcs} \textsuperscript{(2)}, nous avons rencontré un \VUgras{enter}\cc
\VUgras{rement} \textsuperscript{(4)}. En tête du \VUgras{convoi funèbre} marchait le\nl
\VUgras{clergé}, précédant le \VUgras{corbillard} dans lequel était\nl
placé le \VUgras{cercueil}. De nombreux amis accompagnaient\nl
la famille en \VUgras{deuil}.

%%K t14-02-2 p
\VUblancAlinea
Après le passage du cortège, nous avons traversé\nl
la rue devant la grande \VUgras{brasserie} \textsuperscript{(5)}, où de nom\cc
breux \VUgras{consommateurs} buvaient de la \VUgras{bière} sur la\nl
terrasse, à l'ombre de la \VUgras{marquise} \textsuperscript{(6)} vitrée.

%%K t14-03-1 p
\VUblancAlinea
3. --- Jean fut très intrigué par l'\VUgras{écusson} placé\nl
entre le magasin du \VUgras{liquoriste} \textsuperscript{(7)} et celui du \VUgras{mar}\cc
\VUgras{chand de musique} \textsuperscript{(10)}. Il me demanda ce qu'il\nl
signifiait ; je lui répondis qu'il indiquait le \VUgras{consulat}\textsuperscript{(8)}\nl
d'Angleterre, et je lui fis voir le \VUgras{consul} \textsuperscript{(9)}, qui était\nl
sur le balcon.

%%K t14-04-1 p
\VUblancAlinea
4. --- Nous nous sommes arrêtés naturellement\nl
devant l'étalage des \VUgras{instruments} de musique, \VUgras{har}\cc
\VUgras{moniums, orgues} et \VUgras{pianos} ; nous avons regardé\nl
les couvertures illustrées des \VUgras{partitions} et des\nl
\VUgras{morceaux} pour piano, et admiré la \VUgras{lyre} en bois\nl
doré qui sert d'enseigne.

%%K t14-05-1 p
\VUblancAlinea
5. --- Les nuages de poussière soulevés par les\nl
\VUgras{balayeurs} \textsuperscript{(11)} et la \VUgras{balayeuse mécanique} \textsuperscript{(12)}\nl
nous obligèrent à passer rapidement devant les beaux
\parplein
\end{VUpage}

% --- Feuillet 79 (folio 77) -----------------------------------------
\begin{VUpage}[79]{77}
%%K t14-05-1 p suite
\VUcontinue
\VUgras{mobiliers} du \VUgras{marchand de meubles} \textsuperscript{(13)} et devant\nl
l'étalage de \VUgras{poêles} et de \VUgras{lampes} du \VUgras{quincail}\cc
\VUgras{lier} \textsuperscript{(14)}.

%%K t14-06-1 p
\VUblancAlinea
6. --- D'ailleurs, à cet endroit, nous avons été\nl
obligés de quitter le trottoir, car il était encombré par\nl
des colis que l'on descendait d'une \VUgras{voiture de démé}\cc
\VUgras{nagement} \textsuperscript{(16)} pour les porter dans un \VUgras{magasin} \textsuperscript{(15)}\nl
qui venait d'être loué.

%%K t14-06-2 p
\VUblancAlinea
Cela nous empêcha de voir les jolies voitures du\nl
\VUgras{carrossier} \textsuperscript{(17)}, mais non de nous arrêter pour\nl
regarder l'\VUgras{emballeur} \textsuperscript{(19)} faire une caisse pour son\nl
voisin le \VUgras{marchand de cycles et d'automo}\cc
\VUgras{biles} \textsuperscript{(18)} ; Puis, comme il était déjà assez tard, nous\nl
sommes rentrés à la maison.

%%K t14-07-1 p
\VUblancAlinea
7. --- Le lendemain, ma mère proposa à Jean de le\nl
faire photographier et elle me chargea de l'accompa\cc
gner chez le \VUgras{photographe} \textsuperscript{(20)}. Après déjeuner, nous\nl
nous sommes rendus à l'\VUgras{atelier} de l'artiste, où nous\nl
avons été obligés d'attendre assez longtemps. Pour\nl
nous occuper, nous avons regardé les \VUgras{portraits} \textsuperscript{(22)}\nl
pendus au mur.

%%K t14-07-2 p
\VUblancAlinea
Lorsque notre tour est arrivé, l'opération n'a pas\nl
été longue et, dès que mon ami eut fini de poser\nl
devant l'\VUgras{appareil photographique} \textsuperscript{(24)}, nous som\cc
mes redescendus.

%%K t14-08-1 p
\VUblancAlinea
8. --- Au premier étage, nous avons vu par la porte\nl
du \VUgras{coiffeur} \textsuperscript{(23)}, qui était ouverte, son garçon qui\nl
taillait les cheveux d'un client.

%%K t14-08-2 p
\VUblancAlinea
Arrivés dans la rue, nous passâmes devant le\nl
\VUgras{bureau de tabac} \textsuperscript{(24)}. Jean était si amusé par la\nl
\VUgras{lanterne} \textsuperscript{(25)} rouge et la \VUgras{grosse pipe} qui sert\nl
d'\VUgras{enseigne} \textsuperscript{(26)} qu'il alla se heurter contre deux vieux\nl
messieurs qui causaient en prenant une \VUgras{prise} \textsuperscript{(27)}.

%%K t14-09-1 p
\VUblancAlinea
9. --- Les \VUgras{billets de banque} et les \VUgras{pièces de}\nl
\VUgras{monnaie} de tous pays exposés dans la \VUgras{vitrine} du
\parplein
\end{VUpage}

% --- Feuillet 80 (folio 78) -----------------------------------------
\begin{VUpage}[80]{78}
%%K t14-09-1 p suite
\VUcontinue
\VUgras{changeur} \textsuperscript{(29)} attirèrent un moment notre attention ;\nl
mais, comme il était l'heure de goûter, nous sommes\nl
entrés chez le \VUgras{pâtissier} \textsuperscript{(31)}, sans nous arrêter plus\nl
longtemps.

%%K t14-10-1 p
\VUblancAlinea
10. --- A la vue de toutes les friandises que conte\cc
nait le magasin, \VUgras{gâteaux, pains d'épice, biscuits}\nl
et \VUgras{bonbons} de toute sorte, nous ne savions que\nl
choisir. Jean se décida enfin pour les \VUgras{gâteaux} à la\nl
crème, et moi, je ne pris qu'une petite \VUgras{tarte} aux\nl
cerises, car la \VUgras{confiserie} me donne \VUgras{mal aux dents,}\nl
et je ne tiens pas du tout à aller rendre trop souvent\nl
visite au \VUgras{dentiste} \textsuperscript{(30)}.

%%K t14-11-1 p
\VUblancAlinea
11. --- Pour montrer à mon ami une \VUgras{machine à}\nl
\VUgras{écrire}, dont il avait entendu parler, mais qu'il ne con\cc
naissait pas, nous sommes entrés dans le \VUgras{bureau} \textsuperscript{(32)}\nl
de mon \VUgras{cousin} \textsuperscript{(34)}. Nous l'avons trouvé dictant sa\nl
correspondance à son \VUgras{secrétaire} \textsuperscript{(35)}, qui est \VUgras{sténo}\cc
\VUgras{graphe}. A peine une lettre était-elle finie, qu'un\nl
\VUgras{dactylographe} \textsuperscript{(33)} la recopiait rapidement à l'aide\nl
de sa machine. Ne voulant pas interrompre les occu\cc
pations de mon parent, nous ne sommes restés que
\parplein
\end{VUpage}

% --- Feuillet 81 (folio 79) -----------------------------------------
\begin{VUpage}[81]{79}
%%K t14-11-1 p suite
\VUcontinue
quelques instants chez lui.

%%K t14-12-1 p
\VUblancAlinea
12. --- Au-dessous du bureau de mon cousin habite\nl
un grand \VUgras{horloger-bijoutier} \textsuperscript{(37)} qui est en même\nl
temps \VUgras{orfèvre}. Son superbe magasin contient une\nl
quantité d'\VUgras{horloges}, de \VUgras{pendules}, de \VUgras{montres}, de\nl
\VUgras{chaînes} et de \VUgras{bijoux} de prix, tels que \VUgras{colliers} de\nl
perles, \VUgras{bracelets} en or, \VUgras{boucles d'oreilles} et\nl
\VUgras{bagues} ornées de \VUgras{pierres précieuses} et de \VUgras{dia}\cc
\VUgras{mants}. L'une des vitrines est spécialement consa\cc
crée à l'\VUgras{orfèvrerie} et renferme un choix considé\cc
rable de \VUgras{couverts} en argent.

%%K t14-13-1 p
\VUblancAlinea
13. --- En tournant le coin de la rue, je me suis\nl
aperçu que l'on avait changé la \VUgras{plaque} \textsuperscript{(36)} placée\nl
auprès du \VUgras{numéro} de la maison. J'allais en faire\nl
tout haut la réflexion, lorsque les sons joyeux de la\nl
musique militaire se sont fait entendre. Nous nous
\parplein
\end{VUpage}

% --- Feuillet 82 (folio 80) ------------------------------------------
\begin{VUpage}[82]{80}
%%K t14-13-1 p suite
\VUcontinue
sommes mis à courir au devant du régiment, sans\nl
réfléchir qu'il allait passer devant les magasins du\nl
\VUgras{chapelier} \textsuperscript{(39)} et du \VUgras{gantier} \textsuperscript{(40)}, et que nous aurions\nl
été tout aussi bien pour le voir défiler en restant\nl
devant la vitrine de l'\VUgras{opticien} \textsuperscript{(38)}, toute remplie de\nl
\VUgras{lorgnons}, de \VUgras{lunettes} et de \VUgras{lorgnettes}.

%%K t14-14-1 p
\VUblancAlinea
14. --- En tête du \VUgras{régiment} \textsuperscript{(41)} marchait le \VUgras{tam}\cc
\VUgras{bour-major} \textsuperscript{(42)}, suivi par les \VUgras{tambours} \textsuperscript{(43)}, les\nl
\VUgras{clairons} \textsuperscript{(44)} et toute la \VUgras{musique}. Le \VUgras{général} \textsuperscript{(45)},\nl
qui se trouvait sur la place avec son aide-de-camp,\nl
s'était arrêté auprès du \VUgras{jet d'eau} \textsuperscript{(49)} de la \VUgras{fon}\cc
\VUgras{taine} \textsuperscript{(48)}, pour assister au \VUgras{défilé}.

%%K t14-14-2 p
\VUblancAlinea
Les officiers d'\VUgras{état-major} \textsuperscript{(46)}, le \VUgras{colonel} et le\nl
\VUgras{lieutenant-colonel}, le saluèrent de l'épée. Les\nl
\VUgras{commandants} se tenaient à la tête de leur \VUgras{batail}\cc
\VUgras{lon} \textsuperscript{(47)} et chaque \VUgras{capitaine} commandait sa compa\cc
gnie, tandis que les \VUgras{lieutenants, sous-lieutenants}\nl
et \VUgras{sous-officiers} occupaient leur place ordinaire\nl
auprès des hommes.

%%K t14-15-1 p
\VUblancAlinea
15. --- De nombreux passants s'étaient amassés au\nl
\VUgras{carrefour} \textsuperscript{(50)} ; les commerçants, le \VUgras{marchand de}\nl
\VUgras{parapluies} \textsuperscript{(51)}, le \VUgras{teinturier} \textsuperscript{(53)} et l'\VUgras{armurier} \textsuperscript{(54)}\nl
sortirent pour voir les \VUgras{soldats}. Tous les véhicules,\nl
\VUgras{fiacres} \textsuperscript{(52)}, \VUgras{voitures de maître} \textsuperscript{(57)} ainsi que la\nl
\VUgras{tonne d'arrosage} \textsuperscript{(56)}, se rangèrent le long du trot\cc
toir.

%%K t14-15-2 p
Un \VUgras{voyageur de commerce} \textsuperscript{(55)}, qui portait à la\nl
main ses boîtes d'\VUgras{échantillons}, traversa vivement\nl
la rue, et les personnes assises sur l'\VUgras{impériale} \textsuperscript{(59)}\nl
de l'\VUgras{omnibus} \textsuperscript{(58)} se retournèrent pour jouir du\nl
spectacle. Un pauvre \VUgras{chanteur ambulant} \textsuperscript{(60)},\nl
accompagné par un \VUgras{orgue de Barbarie} \textsuperscript{(61)}, dut\nl
interrompre sa \VUgras{chanson}, car le bruit des tambours\nl
couvrait sa voix.

%%K t14-16-1 p
\VUblancAlinea
16. --- Lorsque le régiment arriva devant le\nl
\VUgras{magasin de nouveautés} \textsuperscript{(64)}, les \VUgras{couturières} \textsuperscript{(63)}\nl
et les \VUgras{tailleurs} \textsuperscript{(62)} quittèrent leur \VUgras{machine à cou}\cc
\VUgras{dre} et leur travail pour regarder par la fenêtre. Le
\parplein
\end{VUpage}

% --- Feuillet 83 (folio 81) ------------------------------------------
\begin{VUpage}[83]{81}
%%K t14-16-1 p suite
\VUcontinue
\VUgras{négociant} \textsuperscript{(65)}, qui venait de placer de nouveaux\nl
\VUgras{costumes} \textsuperscript{(66)} dans son \VUgras{étalage}, dut reculer les\nl
pièces de \VUgras{drap} \textsuperscript{(67)} et de \VUgras{toile} installées sur le\nl
trottoir, près de la \VUgras{bouche d'égout} \textsuperscript{(68)}, pour que la\nl
foule ne les renversât pas, et un vieux \VUgras{colporteur} \textsuperscript{(69)},\nl
auquel une concierge marchandait des bas, fut forcé\nl
d'entrer dans un corridor pour terminer sa vente.

%%K t14-16-2 p
\VUblancAlinea
Nous avons accompagné le régiment jusqu'à la\nl
caserne, et, très contents de notre journée, nous\nl
sommes rentrés à l'heure du dîner.

\VUsaut{3.0mm}
\VUfilet{16mm}
\VUsaut{3.0mm}

%%K t14-17-1 p
\VUblancAlinea
17. --- Le jour suivant, mon père nous ayant offert\nl
de nous faire visiter l'imprimerie du journal de la\nl
localité, nous avons accepté avec enthousiasme.

%%K t14-17-2 p
\VUblancAlinea
L'\VUgras{imprimeur} est à la fois \VUgras{libraire} \textsuperscript{(70)}, \VUgras{éditeur},\nl
\VUgras{papetier} et \VUgras{relieur}. Il a installé au premier étage\nl
la \VUgras{rédaction} \textsuperscript{(71)} du journal, ainsi que l'atelier de\nl
\VUgras{gravure}, dans lequel travaillent les \VUgras{lithogra}\cc
\VUgras{phes} \textsuperscript{(72)} et les \VUgras{graveurs} \textsuperscript{(73)} sur cuivre, et l'atelier\nl
de composition, où se trouvent les \VUgras{typographes} \textsuperscript{(74)}.\nl
L'\VUgras{imprimerie} \textsuperscript{(75)} proprement dite, les \VUgras{presses} et\nl
les différentes machines, sont au rez-de-chaussée.

%%K t14-18-1 p
\VUblancAlinea
18. --- En sortant de l'imprimerie, Jean s'arrêta très\nl
intrigué devant la petite boîte vitrée placée sur une\nl
colonne, en face du magasin du \VUgras{mercier} \textsuperscript{(77)} et de\cc
manda à quoi elle servait. Mon père lui expliqua que\nl
c'était un \VUgras{avertisseur d'incendie} \textsuperscript{(76)}. Tout, d'ail\cc
leurs, dans la ville, était pour mon ami un sujet\nl
d'étonnement, depuis la simple \VUgras{borne-fontaine} \textsuperscript{(78)}\nl
et le léger \VUgras{tricycle-porteur} \textsuperscript{(80)} conduit par un\nl
\VUgras{groom} \textsuperscript{(79)} en livrée, jusqu'à la \VUgras{voiture des}\nl
\VUgras{postes} \textsuperscript{(81)}.

%%K t14-19-1 p
\VUblancAlinea
19. --- Pendant que mon père nous quittait un\nl
moment pour parler au \VUgras{percepteur} \textsuperscript{(83)}, qui était\nl
arrêté à causer avec un \VUgras{avocat} \textsuperscript{(82)} et un \VUgras{notaire} \textsuperscript{(84)},\nl
nous nous sommes amusés à suivre des yeux le mou\cc
vement de la rue. Les gens les plus divers passaient\nl
devant nous : un \VUgras{garçon pâtissier} \textsuperscript{(85)}, portant dans\nl
une corbeille un magnifique \VUgras{pâté}, venait derrière un\nl
\VUgras{marchand d'habits} \textsuperscript{(86)} ; un \VUgras{allumeur de réver}\cc
\VUgras{bères} \textsuperscript{(87)} causait avec un \VUgras{gazier} \textsuperscript{(88)} ; une \VUgras{reli}\cc
\VUgras{gieuse} \textsuperscript{(89)}, tenant une petite orpheline par la main,\nl
rentrait à son couvent.

%%K t14-20-1 p
\VUblancAlinea
20. --- Au moment où mon père nous rejoignit,\nl
Jean riait de tout cœur à la vue de la figure bar\cc
bouillée et du bizarre accoutrement de deux \VUgras{ramo}\cc
\VUgras{neurs} \textsuperscript{(91)} tout noirs de suie.

%%K t14-21-1 p
\VUblancAlinea
21. --- Je voulais acheter une plante pour ma mère\nl
à la \VUgras{marchande de fleurs} \textsuperscript{(92)}, mais mon père me\nl
fit remarquer qu'elle serait bien lourde à emporter et\nl
qu'il était préférable de prendre des \VUgras{oranges}. Nous\nl
nous sommes approchés de la \VUgras{marchande} \textsuperscript{(94)} et\nl
lorsque la \VUgras{gouvernante} \textsuperscript{(93)}, qui en choisissait pour\nl
son petit élève, eut fini son achat, nous nous sommes\nl
fait servir.

%%K t14-22-1 p
\VUblancAlinea
22. --- En rentrant à la maison, nous avons ren\cc
contré plusieurs personnes de connaissance : une\nl
vieille amie de ma famille, qui donnait le bras à sa\nl
\VUgras{dame de compagnie} \textsuperscript{(95)}, l'\VUgras{ingénieur} \textsuperscript{(96)} des ponts\nl
et chaussées, et une de mes cousines avec sa \VUgras{bonne}\nl
\VUgras{d'enfant} \textsuperscript{(98)}.

%%K t14-22-2 p
\VUblancAlinea
Puis nous avons croisé la \VUgras{modiste} \textsuperscript{(97)} de ma mère\nl
et deux \VUgras{moines} \textsuperscript{(99)} étrangers à la ville, qui nous\nl
abordèrent pour demander à mon père le chemin de\nl
la gare.

\VUsaut{6.0mm}
\VUfilet{22mm}
\end{VUpage}
